\documentclass[10pt, xcolor={dvipsnames}]{beamer}

%========================
% CẤU HÌNH HỆ THỐNG
%========================
\usepackage[utf8]{inputenc}
\usepackage[T5]{fontenc}
\usepackage[vietnamese]{babel}
\usepackage{tikz}
\usetikzlibrary{arrows.meta, positioning, calc, shapes.geometric, shadows, patterns}
\usepackage{amsmath, amssymb, amsthm}
\usepackage{tcolorbox}
\tcbuselibrary{skins, shadows, themes}

%========================
% TÔNG MÀU GỐC XANH - CAM
%========================
\usetheme{Madrid}
\definecolor{mainblue}{RGB}{0, 102, 104} 
\definecolor{lightorange}{RGB}{255, 140, 0} 
\definecolor{correctgreen}{RGB}{0, 153, 51}

\setbeamercolor{palette primary}{bg=mainblue, fg=white}
\setbeamercolor{palette secondary}{bg=lightorange, fg=white}
\setbeamercolor{structure}{fg=mainblue}

% Hộp nội dung đặc sắc
\newtcolorbox{gamebox}[2]{
    enhanced, colback=white, colframe=#1, title=#2,
    fonttitle=\bfseries\large, arc=3pt, boxrule=1pt, 
    shadow={1mm}{-1mm}{0mm}{black!10}
}

\begin{document}

%--- SLIDE MỞ ĐẦU ---
\begin{frame}[plain]
    \centering
    \begin{tikzpicture}
        \node[starburst, fill=lightorange, draw=mainblue, line width=2pt, minimum size=4cm] at (0,0) {\bfseries \Huge \color{white} GAME};
    \end{tikzpicture}
    \vspace{0.4cm}
    \begin{gamebox}{mainblue}{NHÓM THỰC HIỆN}
        \centering \Large \textbf{Tung Tung Tung Sahur} \\
        \vspace{0.2cm}
        \small 10 Câu Hỏi Thử Thách Đồ Thị Thú Vị
    \end{gamebox}
\end{frame}

%--- CÂU 1 ---
\begin{frame}{Câu 1: Bài toán Bắt tay}
    \begin{gamebox}{mainblue}{Câu hỏi}
        Trong một bữa tiệc có 10 người. Nếu mỗi người đều bắt tay với tất cả những người còn lại đúng 1 lần, có tổng cộng bao nhiêu cái bắt tay?
    \end{gamebox}
    \begin{itemize}
        \item[\textbf{A.}] 100 \quad \quad \quad \quad \textbf{B.} 90
        \item[\textbf{C.}] \only<1>{45}\only<2>{\textbf{\color{correctgreen}45 (Đúng!)}} \quad \quad \textbf{D.} 20
    \end{itemize}
    \only<2>{\textit{\color{mainblue}Giải thích: Đây là đồ thị đầy đủ $K_{10}$, số cạnh là $C^{2}_{10} = 45$.}}
\end{frame}

%--- CÂU 2 ---
\begin{frame}{Câu 2: Bảy cây cầu Königsberg}
    \begin{gamebox}{lightorange}{Câu hỏi}
        Leonhard Euler đã chứng minh không thể đi qua 7 cây cầu mà mỗi cầu chỉ đi 1 lần. Lý do chính là gì?
    \end{gamebox}
    \begin{itemize}
        \item[\textbf{A.}] Đồ thị không liên thông.
        \item[\textbf{B.}] \only<1>{Có quá nhiều đỉnh bậc lẻ.}\only<2>{\textbf{\color{correctgreen}Có quá nhiều đỉnh bậc lẻ (4 đỉnh).}}
        \item[\textbf{C.}] Đồ thị là một cây.
        \item[\textbf{D.}] Có chu trình Euler.
    \end{itemize}
\end{frame}

%--- CÂU 3 ---
\begin{frame}{Câu 3: Bài toán Tô màu bản đồ}
    \begin{gamebox}{mainblue}{Câu hỏi}
        Số màu tối thiểu cần thiết để tô màu bất kỳ bản đồ phẳng nào sao cho các vùng giáp ranh không cùng màu là bao nhiêu?
    \end{gamebox}
    \begin{itemize}
        \item[\textbf{A.}] 2 màu \quad \quad \quad \quad \textbf{B.} 3 màu
        \item[\textbf{C.}] \only<1>{4 màu}\only<2>{\textbf{\color{correctgreen}4 màu (Định lý 4 màu)}} \quad \textbf{D.} 5 màu
    \end{itemize}
\end{frame}

%--- CÂU 4 ---
\begin{frame}{Câu 4: Mạng xã hội}
    \begin{gamebox}{lightorange}{Câu hỏi}
        Trong đồ thị mạng xã hội (Facebook), quan hệ "Bạn bè" thuộc loại đồ thị nào?
    \end{gamebox}
    \begin{itemize}
        \item[\textbf{A.}] Có hướng.
        \item[\textbf{B.}] \only<1>{Vô hướng.}\only<2>{\textbf{\color{correctgreen}Vô hướng (Kết bạn là 2 chiều).}}
        \item[\textbf{C.}] Đồ thị rỗng.
        \item[\textbf{D.}] Đồ thị phân đôi hoàn toàn.
    \end{itemize}
\end{frame}

%--- CÂU 5 ---
\begin{frame}{Câu 5: Bài toán Con mã trên bàn cờ}
    \begin{gamebox}{mainblue}{Câu hỏi}
        Lộ trình của con mã đi qua tất cả các ô bàn cờ đúng 1 lần tương ứng với khái niệm nào?
    \end{gamebox}
    \begin{itemize}
        \item[\textbf{A.}] Đường đi Euler.
        \item[\textbf{B.}] \only<1>{Đường đi Hamilton.}\only<2>{\textbf{\color{correctgreen}Đường đi Hamilton.}}
        \item[\textbf{C.}] Cây khung nhỏ nhất.
        \item[\textbf{D.}] Luồng cực đại.
    \end{itemize}
\end{frame}

%--- CÂU 6 ---
\begin{frame}{Câu 6: Bậc của Đỉnh}
    \begin{gamebox}{lightorange}{Câu hỏi}
        Trong một nhóm bạn, có thể tồn tại một nhóm mà số người có số lượng bạn bè là số lẻ lại là một số lẻ không?
    \end{gamebox}
    \begin{itemize}
        \item[\textbf{A.}] Có thể.
        \item[\textbf{B.}] \only<1>{Không thể.}\only<2>{\textbf{\color{correctgreen}Không thể (Định lý bắt tay).}}
        \item[\textbf{C.}] Tùy vào số lượng người.
        \item[\textbf{D.}] Chỉ khi đồ thị phẳng.
    \end{itemize}
\end{frame}

%--- CÂU 7 ---
\begin{frame}{Câu 7: Độ lệch Google Maps}
    \begin{gamebox}{mainblue}{Câu hỏi}
        Thuật toán phổ biến nhất mà Google Maps dùng để tìm đường ngắn nhất là gì?
    \end{gamebox}
    \begin{itemize}
        \item[\textbf{A.}] Kruskal \quad \quad \quad \quad \textbf{B.} Prim
        \item[\textbf{C.}] \only<1>{Dijkstra}\only<2>{\textbf{\color{correctgreen}Dijkstra (và biến thể A*)}} \quad \textbf{D.} Ford-Fulkerson
    \end{itemize}
\end{frame}

%--- CÂU 8 ---
\begin{frame}{Câu 8: Cây Phả hệ}
    \begin{gamebox}{lightorange}{Câu hỏi}
        Cây phả hệ của một dòng họ (không có hôn nhân cận huyết) có chứa chu trình không?
    \end{gamebox}
    \begin{itemize}
        \item[\textbf{A.}] Có.
        \item[\textbf{B.}] \only<1>{Không.}\only<2>{\textbf{\color{correctgreen}Không (Cây không có chu trình).}}
        \item[\textbf{C.}] Chỉ ở đời thứ 3.
        \item[\textbf{D.}] Có, nếu tính cả họ hàng xa.
    \end{itemize}
\end{frame}

%--- CÂU 9 ---
\begin{frame}{Câu 9: Bố trí máy chủ}
    \begin{gamebox}{mainblue}{Câu hỏi}
        Để nối 5 máy chủ sao cho máy nào cũng thông với máy nào bằng ít dây cáp nhất, cần bao nhiêu dây?
    \end{gamebox}
    \begin{itemize}
        \item[\textbf{A.}] 5 dây \quad \quad \quad \quad \textbf{B.} 10 dây
        \item[\textbf{C.}] \only<1>{4 dây}\only<2>{\textbf{\color{correctgreen}4 dây (Đặc điểm của Cây)}} \quad \textbf{D.} 6 dây
    \end{itemize}
\end{frame}

%--- CÂU 10 ---
\begin{frame}{Câu 10: Người xa lạ}
    \begin{gamebox}{lightorange}{Câu hỏi}
        Giả thuyết "6 chặng phân cách" (Six Degrees of Separation) nói rằng bất kỳ ai cũng kết nối với nhau qua tối đa bao nhiêu người trung gian?
    \end{gamebox}
    \begin{itemize}
        \item[\textbf{A.}] 3 người \quad \quad \quad \quad \textbf{B.} \only<1>{6 người}\only<2>{\textbf{\color{correctgreen}6 người}}
        \item[\textbf{C.}] 10 người \quad \quad \quad \quad \textbf{D.} 100 người
    \end{itemize}
\end{frame}

%--- SLIDE KẾT THÚC ---
\begin{frame}[plain]
    \centering
    \begin{tikzpicture}
        \node[circle, shading=ball, ball color=lightorange, minimum size=3cm] at (0,0) {\Huge \bfseries \color{white} THANKS!};
    \end{tikzpicture}
    
    \vspace{0.6cm}
    \begin{gamebox}{mainblue}{LỜI CẢM ƠN}
        \centering
        Cảm ơn thầy và các bạn đã dành thời gian theo dõi \\ phần trình bày của nhóm chúng em!
    \end{gamebox}
    
    \vspace{0.4cm}
    \textbf{Nhóm sinh viên thực hiện:} \\
    \textcolor{mainblue}{\Large \textbf{Tung Tung Tung Sahur}} \\
    \vspace{0.4cm}
    \textit{"Đồ thị là cuộc sống, hãy kết nối chúng lại!"}
\end{frame}

\end{document}